\documentclass{llncs}

\usepackage{amsmath,amssymb,amsfonts}
\usepackage{mathtools}
\usepackage{tikz}
\usepackage{enumitem}
\usepackage{algorithm}

\begin{document}
\title{Zeeperio: Verifying Elections using Poly-IOP Gadgets and On-Chain Verification}
\titlerunning{Zeeperio}



\keywords{verifiable voting, zero-knowledge proofs, SNARKs, polynomial interactive oracle proofs, KZG commitments, blockchain}

\begin{abstract}
\end{abstract}


\section{Introduction}

The legitimacy of any democratic election relies on its governing structure to provide trust, which is done by observing the counting process. 
As noted by \_\_, the transparency of the counting process is not just a preference, rather a requirement: \emph{"The counting of ballots
and determination of the results of voting must be a transparent process that is open to observation"}. Without independent scrutiny, the voters cannnot maintain confidence in the security of their ballots, or the accuracy of the tallies.

Despite these clear requirements, modern election administrations have moved away from observability. 
In jurisdictions such as Canada, municipal and provincial elections increasingly rely on \emph{optical scanners} for paper ballots and \emph{remote online voting systems}. 

In both scenarios, the counting process is performed by proprietary software running on closed hardware. 
A voter casting a ballot online, or a voter feeding a paper ballot into a tabulator both face the same problem: 
They must blindly trust the vendor's software is free of any bugs, has not been tempered with, and the counting process is done accurately.
This \emph{black box} nature means individual voters have little to no ability to independently verify their own ballot was recorded and tallied correctly.


\section{Background and Related Work}

\subsection{E2E verifiable voting}

End-to-end verifiable (E2E) voting protcols were developed to address this gap by providing the following properties.
\begin{description}
  \item \textbf{Cast-as-intended:} The voter can verify their vote was correctly captured by the system.
  \item \textbf{Recorded-as-cast:} The voter can verify their vote was included in the final tally.
  \item \textbf{Tallied-as-recorded:} Any observer can verify the announced outcome is consistent w.r.t. the ballots.
\end{description}

These voting systems typically offer a public \emph{bulletin board} where all the cryptographic artifacts from the election are posted.
Voters then obtain some form of \emph{receipt} that allows them to locate their ballot data on the board and check that it has not been altered.
Finally, observers can use proofs or audits to check that the final tally is consistent with the posted artifacts.

Paper-based systems such as Pr{\^e}t-\`a-Voter and Scantegrity add additional constraints to the paper ballot by adding random candidate permutations and confirmation codes.
These schemes rely on mix-nets, homomorphic tallying and zero-knowledge proofs to show the final tally corresponds to the posted data, without revealing how any individual voted.

Online systems such as Helios encrypt each ballot under a public key and use homomorphic encryption for tallying.
Voters receive a hash for their encrypted ballot, which they can later check on the bulletin board.
Helios also uses mix-nets and decryption proofs to ensure the outcome matches the encrypted ballots.

Civitas and other coercion-resistant systems add further layers of cryptographic security to allow voters to deny coercion. These schemes provide strong guarantees but are complex to implement and verify in practice.

\subsection{Eperio}

Eperio is an audit-focused protocol that aims to minimize the cryptographic verification on auditors.
Instead of heavy zero-knowledge proofs, it organizes election data into collections of files that can be checked with simple cryptography and spreadsheets.
Zeeperio follows Eperio's philosophy of keeping verification simple for the verifier, but uses succinct proofs instead of spreadsheets.

\subsection{Succinct proofs and Poly-IOPs}

Modern SNARK systems such as PLONK, Groth16, and Marlin are all based on the Polynomial Interactive Oracle Proof (Poly-IOP) model.
In a Poly-IOP, the prover commits to low-degree polynomials representing the witness. Then, the verifier issues queries to check the resulting evaluations satisfy a system of polynomial identities.
Polynomial commitment schemes allow these evaluations to be proven succinctly.

\subsection{Zeeperio: A succinct, On-Chain Approach}
Zeeperio shifts from statistical audits to cryptographic arguments of knowledge. 
At a high level, Zeeperio uses the Poly-IOP model to encodes the entire state of an election as a small set of polynomials. 
Instead of writing a generic arithmetic circuit, Zeeperio works directly at the polynomial level to design custom constraints for ballot structure, sums over ballots, and tallies.
Then, using the KZG commitment scheme, Zeeperio allows the Election Authority (EA) to generate a single, succinct proof that the election outcome is correct.
This proof is constant in size and verification time, regardless of how many voters cast their ballots.

This efficiency allows for the verification logic to be encoded into an
Ethereum smart contract. Instead of requiring observers to download data and
run software, the \emph{public audit} becomes a transaction on a the blockchain. 
If the smart contract accepts the proof, the election tally is mathematically guaranteed to be consistent with the committed ballots. 
This lowers the barrier to verification from running complex cryptographic software to checking a simple status flag on the Ethereum block explorer.

\section{Threat Model and Contributions}

Before detailing the protocol mechanics, we formalize the operational environment and the specific contributions of this work relative to the state of the art.

\subsection{The RealCryptik Threat Model}
Zeeprio adopts the \textit{RealCryptik} threat model. 
Traditional E2E literature often models the EA as a powerful adversary that must be prevented from learning individual votes at all costs. 
However, this model often ignores the reality of elections. In many scenarios, the EA has the full capability to deanonymize votes.

The RealCryptik model explicitly assumes that the EA \textbf{has access to who voted for whom}. Under this assumption, the primary security goal shifts from absolute privacy to absolute \textbf{integrity}. 
The system must guarantee that even an EA with full view of voter choices cannot manipulate an election without detection. This model reflects the reality of \emph{black box} vendors while enforcing the verification constraints necessary for public trust. 
While public privacy is maintained through zero-knowledge proofs, the protocol also prevents the EA from manipulating the outcome based on their knowledge.

\subsection{Contributions}
This work presents the following contributions.

\begin{description} 
\item \textbf{Succinct universal verification:} All election constraints are encoded as a single Poly-IOP relation and proven via a zkSNARK proof, independent of the number of voters.
\item \textbf{Enhanced privacy and dispute resolution:}
  Confirmation codes are committed to a single polynomial and are only selectively opened when needed.
  After a voter has cast their ballot, they receives a receipt containing a confirmation code to which they can check it appears in the public bulletin board. 
 This is done in a way that prevents the receipt from being used as a proof of how the voter voted.
  \item \textbf{On-chain verification:}
  The final proof is verified by an Ethereum smart contract. This removes the need for a proprietary verifier software and anchors the election result in a public, tamper-evident ledger.
\end{description}


\section{Preliminaries}

We now introduce the cryptographic tools used in Zeeperio.

\subsection{Polynomial interactive oracle proofs}

A polynomial interactive oracle proof (Poly-IOP) is an interactive proof where the prover sends commitments to polynomials and the verifier has oracle access to these polynomials through evaluation queries.
Formally, a Poly-IOP protocol for a relation $\mathcal{R}$ between public input $x$ and witness $w$ follows:
\begin{enumerate}[label=(\roman*)]
  \item The prover, given $x$ and $w$, computes polynomials $f_1,\dots,f_m$ and sends commitments to them.
  \item The verifier sends random challenges $r_1,\dots,r_\ell$.
  \item The prover answers with evaluations $f(r_\ell)$ and proofs of correct evaluation.
  \item The verifier checks the evaluations satisfy a system of polynomial identities and the evaluation proofs are valid.
\end{enumerate}
If all checks pass, the verifier accepts the statement \emph{there exists $w$ such that $(x,w) \in \mathcal{R}$}.

\subsection{KZG polynomial commitments}

Zeeperio uses the KZG (Kate-Zaverucha-Goldberg) polynomial commitment scheme.
We recall its interface at a high level.
A trusted setup is used to produces a structured reference string (SRS) consisting of group elements
\[
  \{g, g^\tau, g^{\tau^2}, g^{\tau^d}\} \in G_1, \qquad \{h, h^\tau\} \in G_2,
\]
for a secret $\tau$.

Given a polynomial $f(X) = \sum_{i=0}^d f_i X^i$, the commitment is
\[
  \mathsf{Com}(f) = \sum_{i=0}^d f_i [\tau^i] \in G_1.
\]
To open $\mathsf{Com}(f)$ at a point $z \in \mathbb{F}$, the prover computes $y = f(z)$ and the quotient
\[
  q(X) = \frac{f(X) - y}{X - z}
\]
and returns an opening proof $\pi = \mathsf{Com}(q)$.
The verifier checks that
\[
  e(\mathsf{Com}(f) - y, h) = e(\pi, [h^\tau - z]),
\]
where $e : G_1 \times G_2 \to G_T$ is a bilinear pairing.

In Zeeperio we use KZG both to commit to the election polynomials.
We combine multiple constraints into a single quotient polynomial using random coefficients and prove that it vanishes on $\Omega$ by checking a single evaluation.

\subsection{Zero-knowledge and Fiat--Shamir}

We can apply the Fiat-Shamir heuristic to the Poly-IOP to make the protcol non-interactive.
This can be done by replacing the verifier's random challenges by outputs of a hash function on the transcript.
To achieve zero-knowledge, we add random blinding polynomials to all witness polynomials before committing, and we ensure that any information opened to the verifier is public.
Therefore, the only meaningful information revealed by the proof is that the constraints are satisfied, while keeping the underlying ballot structure and confirmation codes hidden.

\subsection{Ethereum verification model}

Finally, we briefly outline the on-chain verification model.
Ethereum provides precompiled contracts for the BN254 pairing and scalar multiplication.
The verification contract uses the SRS and the structure of the combined pairing equation corresponding to the KZG openings.
Given the public input and the proof, the contract performs a simple pairing check.
If the check passes, the election result has been cryptographically verified.

\section{The Zeeperio Protocol}

We now describe the Zeeperio protocol in detail, focusing on the data representation and the polynomial gadgets used to enforce election integrity.

\subsection{Election Data Representation}
We consider an election with $\mathcal{B}$ ballots and $\mathcal{C}$ candidates. There are $\mathsf{P} = \mathcal{B} \cdot \mathcal{C}$ unique ballot-candidate positions listed in canonical order. Let $j \in \{0, \dots, \mathcal{B}-1\}$ denote the ballot index and $k \in \{0, \dots, \mathcal{C}-1\}$ the candidate index within a ballot. Then $i = j \cdot \mathcal{C} + k$ iterates over each ballot-candidate position $P_i$.

Each position is associated with four data columns:
\begin{enumerate}
    \item \textbf{Ballot Column ($\mathsf{B}_\text{ID}$):} Identifies the unique ballot $\mathcal{B}_j$.
    \item \textbf{Confirmation Code Column ($\mathsf{C}_\text{confirm}$):} A unique confirmation code issued to each position. The voter receives the code corresponding to their selection.
    \item \textbf{Mark Column ($\mathsf{M}$):} Indicates if the position was selected as the cast vote.
    \item \textbf{Audit Column ($\mathsf{A}$):} Indicates if the ballot $\mathcal{B}_j$ was selected for audit.
\end{enumerate}

For efficient Fast Fourier Transform (FFT) operations, we require an evaluation domain size $n$ that is a power of two. 
We set $n = 2^t$ where $t = \lceil \log_2 \mathsf{P} \rceil$. 
We pad the columns with zeros for indices $i$ where $\mathsf{P} \le i < n$. 
The columns are interpolated into polynomials $\mathsf{B}_\text{ID}(X), \mathsf{C}_\text{confirm}(X), \mathsf{M}(X), \mathsf{A}(X)$ over the domain $\Omega = \{1, \omega, \dots, \omega^{n-1}\}$ where $\omega$ is a primitive $n$th root of unity in $\mathbb{F}_q$.

\subsection{Setup}

\begin{definition}[q-strong Diffie--Hellman (q-sDH)]
Given a tuple 
\(
(g, g^x, g^{x^2}, \ldots, g^{x^q}) \in \mathbb{G}^{q+1}
\),
where \(x \gets_{\$} \mathbb{Z}_q\), q-sDH expects a pair
\(
(g^{\frac{1}{x+c}}, c) \in \mathbb{G} \times \mathbb{Z}_q
\).
q-sDH is believed to be intractable, that is, for any PPT adversary
\(\mathcal{A}\), it is hard to find such a \(c\). Formally,
\begin{equation}
  \Pr\bigl[
    O = g^{\frac{1}{x+c}}
    \,\bigm|\,
    (O,c) \gets \mathcal{A}(g, g^x, g^{x^2}, \ldots, g^{x^q})
  \bigr]
  \leq \mathsf{negl}(\lambda).
  \label{eq:q-sdh}
\end{equation}
\end{definition}

In the setup phase, the election authority (EA) performs the following steps.
\begin{enumerate}[label=(S\arabic*)]
  \item Obtains KZG SRS from publicly run Ethereum Powers-of-Tau ceremony.
  \item Fixes the finite field $\mathbb{F}_q$, root-of-unity domain $\Omega$ of size $n$, and encoding of election data into polynomials as in Section~3.
  \item Generates a set of unique confirmation codes and assigns each $\mathsf{C}_confirm$ code to each position on the ballot.
  \item Deploys an Ethereum smart contract that uses the SRS and the verification logic for the Zeeperio proof.
\end{enumerate}

\subsection{Voting phase}

In the voting phase, voters receive a paper ballot $B_j$ with a unique serial number.  Voters mark the position corresponding to their chosen candidate and cast their ballots. They receive the confirmation code associated with the candidate they chose.
At the end of the voting phase, the EA holds the full table of positions with their ballot IDs, audit flags, mark bits, and confirmation codes.


\subsection{Post-election}

After the polls close, the EA constructs the polynomials $\mathsf{A}(X)$, $\mathsf{M}(X)$, $\mathsf{B}_\mathsf{ID}(X)$, $\mathsf{C}_{\mathsf{confirm}}(X)$ and produces a succinct proof that satisfies the election constraints.

\subsection{Constraint Enforcement with Vanishing Polynomials}
We enforce constraints over the domain $\Omega$ using a vanishing polynomial $Z_\Omega(X) = X^n - 1$. 
A constraint $C(X) = 0$ holdsover $\Omega$ iff $C(X)$ is divisible by $Z_\Omega(X)$. 
In the Poly-IOP model, the prover computes the quotient polynomial $Q(X) = C(X) / Z_\Omega(X)$, commits to it, and the verifier checks the identity at a random point.
For constraints that apply only to specific subsets of the domain, we use specialized vanishing polynomials. 

\subsection{Audit Column Constraints}
The Audit Column $\mathsf{A}$ must satisfy the following properties to ensure the integrity of the audit process.

\subsubsection{Padding Consistency}
We must ensure that the padded values in the Audit Column are zero. We define a constraint specific vanishing polynomial $Z(X)$ with roots at the valid indices $\{\omega^0, \dots, \omega^{\mathsf{P}-1}\}$:
\begin{equation}
    Z(X) = \prod_{i=0}^{\mathsf{P}-1} (X - \omega^i)
\end{equation}
The constraint requires $\mathsf{A}(X)$ is zero wherever $Z(X)$ is non-zero. This is enforced via the domain vanishing polynomial:
\begin{equation}
    \mathsf{P}_\text{vanish1}(X) = \mathsf{A}(X) \cdot Z(X) = 0 \quad \forall X \in \Omega
\end{equation}

\subsubsection{Binary Values}
Every entry in $\mathsf{A}$ must be binary (0 or 1). This is enforced via the domain vanishing polynomial:
\begin{equation}
    \mathsf{P}_\text{vanish2}(X) = \mathsf{A}(X) \cdot (\mathsf{A}(X) - 1) \quad \forall X \in \Omega
\end{equation}
\subsection{Mark Column Constraints}


\section{Security Analysis}

\subsection{Completeness}

\begin{theorem}[Completeness]
    If an honest Prover $\mathcal{P}$ computes the accumulators $Acc_{T_1}(X)$ and $Acc_{T_2}(X)$ according to the backward accumulation rule defined in the protocol, then the Verifier $\mathcal{V}$ accepts with probability 1.
\end{theorem}

\begin{proof}
    The honest prover constructs the backward accumulator such that $Acc_{T_1}(\omega^i) = \prod_{j=i}^{n-1} T_1(\omega^j)$. The protocol enforces constraints via a batched vanishing polynomial $P_{vanish}(X)$.
    \begin{enumerate}
        \item \textbf{Boundary Consistency:} At the last element $X = \omega^{n-1}$, the accumulator is defined as $Acc_{T_1}(\omega^{n-1}) = T_1(\omega^{n-1})$. Consequently, the term $(Acc_{T_1}(X) - T_1(X))$ is zero at this point, satisfying the boundary constraint $P_{vanish1}$.
        \item \textbf{Recursive Consistency:} For all $X \in \Omega$ except $X=\omega^{n-1}$, the recursive definition holds: $Acc_{T_1}(\omega^i) = T_1(\omega^i) \cdot Acc_{T_1}(\omega^{i+1})$. This ensures the constraint $P_{vanish3}$ evaluates to zero on the relevant domain.
        \item \textbf{Product Equality:} Since the tuples are permutations, the total products over $\Omega$ are equal. Thus, $Acc_{T_1}(1) = \prod_{x \in \Omega} T_1(x) = \prod_{x \in \Omega} T_2(x) = Acc_{T_2}(1)$, satisfying the equality constraint $P_{vanish5}$.
    \end{enumerate}
    Since all individual constraint polynomials vanish on $\Omega$, $P_{vanish}(X)$ is divisible by the vanishing polynomial $Z_H(X) = X^n - 1$. Therefore, the quotient check $P_{vanish}(\zeta) - Q(\zeta)Z_H(\zeta) \stackrel{?}{=} 0$ holds for any $\zeta \in \mathbb{F}_q$.
    \qed
\end{proof}

\subsection{Knowledge Soundness}

\begin{theorem}[Knowledge Soundness]
    The protocol satisfies knowledge soundness in the Algebraic Group Model. If a malicious prover ${\mathcal{P}}$ convinces $\mathcal{V}$ with non-negligible probability $\epsilon$, there exists an efficient Extractor $\mathcal{E}$ that outputs polynomials $T(X)$ and $T'(X)$ such that the set of their evaluations on $\Omega$ are permutations of each other.
\end{theorem}

\begin{proof}
    In the AGM, the Extractor $\mathcal{E}$ observes the commitments sent by ${\mathcal{P}}$. Since ${\mathcal{P}}$ outputs valid commitments $[Acc_{T_1}], [Acc_{T_2}]$, it must know the underlying polynomials.
    
    If $\mathcal{V}$ accepts, the equation $P_{vanish}(\zeta) = 0$ holds at a random challenge $\zeta \leftarrow \mathbb{F}_q$. By the Schwartz-Zippel Lemma, this implies that $P_{vanish}(X)$ is divisible by $Z_H(X)$ with overwhelming probability. 
    We unroll these constraints to prove the validity of the statement:

    \begin{enumerate}
        \item \textbf{Unrolling the Accumulator:}
        From the boundary constraint ($P_{vanish1} \equiv 0$), we have $Acc_{T_1}(\omega^{n-1}) = T_1(\omega^{n-1})$.
        From the constraint ($P_{vanish3} \equiv 0$), for any $x \in \Omega$ except $X= \omega^{n-1}$, $Acc_{T_1}(x) = T_1(x) \cdot Acc_{T_1}(\omega x)$.
        With this backwards accumulation, we determine that $Acc_{T_1}(1)$ stores the complete product:
        \begin{equation}
            Acc_{T_1}(1) = \prod_{i=0}^{n-1} T_1(\omega^i)
        \end{equation}
        Similarly, we derive $Acc_{T_2}(1) = \prod_{i=0}^{n-1} T_2(\omega^i)$.

        \item \textbf{Equality of Multisets:}
        The final constraint ($P_{vanish5} \equiv 0$) ensures $Acc_{T_1}(1) = Acc_{T_2}(1)$. Substituting the definitions of $T_1$ and $T_2$:
        \begin{equation}
            \prod_{x \in \Omega} (r - T(x)) = \prod_{x \in \Omega} (r - T'(x))
        \end{equation}
        This equation implies equality between two polynomials in the variable $r$. Since $r$ is chosen uniformly at random by the verifier \emph{after} the commitments to $T(X)$ and $T'(X)$, the equality holds identically with overwhelming probability. 
        Therefore, these polynomials are considered to be \emph{set equal}. Thus, $T'$ is a valid permutation of $T$.
    \end{enumerate}
    \qed
\end{proof}

\subsection{Honest-Verifier Zero-Knowledge}

\begin{theorem}[HVZK]
    The protocol is Honest-Verifier Zero-Knowledge. There exists a Simulator $\mathcal{S}$ that outputs a transcript indistinguishable from a real prover's trace.
\end{theorem}

\begin{proof}
   The Simulator $\mathcal{S}$ operates as follows:

\begin{enumerate}
    \item \textbf{Challenge Sampling:} $\mathcal{S}$ samples random challenges $r, \beta, \zeta \leftarrow \mathbb{F}_q$ exactly as the honest verifier would.
    
    \item \textbf{Simulate Commitments:} $\mathcal{S}$ outputs random group elements for $[Acc_{T_1}], [Acc_{T_2}], [T_1], [T_2]$. Due to the hiding property of KZG (with blinding), these are indistinguishable from real commitments.
    
    \item \textbf{Simulate Evaluations:}
    \begin{itemize}
        \item $\mathcal{S}$ chooses random values for evaluations $\bar{T}_1, \bar{T}_2, \overline{Acc}_{T_1}, \overline{Acc}_{T_2}$ at $\zeta$ and required shifted points (e.g., $\omega\zeta$). This is valid because the blinded polynomials act as One-Time Pads at any point $\zeta \notin \Omega$.
        
        \item $\mathcal{S}$ must ensure the verification equation holds: $P_{vanish}(\zeta) - Q(\zeta)Z_H(\zeta) = 0$.
        
        \item $\mathcal{S}$ computes the value of $P_{vanish}(\zeta)$ using the randomly sampled evaluations and the formulas for $Z_1, Z_2, Z_3$ evaluated at $\zeta$.
        
        \item $\mathcal{S}$ then \textit{sets} the evaluation of the quotient polynomial $\bar{Q} = Q(\zeta)$ specifically to satisfy the equation:
        \[
        \bar{Q} = \frac{P_{vanish}(\zeta)}{Z_H(\zeta)}
        \]
        Since $\zeta$ is random, $Z_H(\zeta) \neq 0$, making this division valid.
    \end{itemize}
    
    \item \textbf{Transcript Consistency:} The tuple of commitments, challenges, and evaluations $(\dots, \bar{Q})$ satisfies the algebraic verification check by construction. Because the evaluations of blinded polynomials are uniformly distributed, the simulated transcript is statistically indistinguishable from a real transcript. 
\end{enumerate}
\end{proof}

\section{Implementation}

\section{Future Work}

We conclude with several directions for extending and deploying Zeeperio.

\paragraph{Integration with optical scanners}

\paragraph{Online voting}

\end{document}
